\documentclass[10pt,a4paper]{amsart}
\usepackage[margin=19mm]{geometry}
\usepackage{amsmath,amssymb,mathtools}
\usepackage[scr=rsfs,cal=euler]{mathalfa}
\usepackage{xfrac}
\usepackage[unicode,hidelinks,bookmarks=false]{hyperref}
\newcommand{\ie}{i.e.\ }
\newcommand{\cf}{cf.\ }
\newcommand{\resp}{respectively\ }
\newcommand{\Subsection}{\subsection}
\providecommand{\nobreakdash}{\nobreak\hbox{-}}
\setlength{\emergencystretch}{2em}
\multlinegap0pt
\usepackage{bm}
\usepackage{bbm}
\usepackage{dsfont}
\usepackage{xspace}
\usepackage{cancel}
\usepackage{mathtools}
\usepackage{amsthm}
\makeatletter
\@ifundefined{theorem}{\newtheorem{theorem}{Theorem}}{}
\@ifundefined{proposition}{\newtheorem{proposition}[theorem]{Proposition}}{}
\makeatother
\newcommand{\FF}{\mathcal F}
\newcommand{\eps}{\varepsilon}
\title[Fibonacci, Dirichlet, and Gauss in a single sum]{Fibonacci, Dirichlet, and Gauss in a single sum}
\author{Benoit Cloitre}
\date{Source: arXiv:2607.20960v1 (CC BY 4.0)}
\begin{document}
\maketitle

\noindent\emph{Source and licence.} Fibonacci, Dirichlet, and Gauss in a single sum. Version arXiv:2607.20960v1 (CC BY 4.0), distributed under the Creative Commons Attribution 4.0 International licence, as recorded in the arXiv licence metadata for this exact version and shown on the abstract page https://arxiv.org/abs/2607.20960v1. Full rights notice: licenses/arxiv-2607.20960-CC-BY-4.0.txt.\par\smallskip

\noindent\textit{Editorial scope.} Counted unit: the Proposition whose source label is \texttt{prop:converse} in the article (source0.tex lines 505--519, source label \texttt{prop:converse}), delivered together with the article's own complete original proof of it (source0.tex lines 521--590, one \texttt{proof} environment of 70 lines). No mathematics is written, repaired or re-derived: statement and proof are the article's own text line for line. Required context retained uncounted: eq:def-hyperbola (lines 101--106); eq:odd-main (lines 152--156); eq:even-main (lines 158--164). Every definition, display and label of those lines is retained unchanged. Omitted from this package and not redistributed: 1--100, 107--151, 157--157, 165--504, 520--520, 591--899. Nothing inside a retained excerpt is omitted or altered: every retained body is delivered exactly as the article's lines are written, so no omission receipt of any kind accompanies this package. Citations: the retained excerpts carry no citation command at all, so no bibliography entry of the article is transcribed and no reference is redistributed. Reference adaptation: none was needed and none was made; every reference inside the delivered unit resolves inside it, so no unresolved reference or citation is produced. Printed slips kept literal: no printed word, symbol, hypothesis or number of the counted statement or of its proof was altered. Layout and class: the article's own class is \texttt{article} with packages including fontenc, lmodern, amsmath, amssymb, amsthm, microtype, geometry, hyperref; the wrapper is the lane's pinned \texttt{amsart} editorial wrapper at 10pt on A4, which restates the theorem-like environments the retained text needs and adds the article's own macro definitions that the retained text needs (\texttt{\textbackslash FF}, \texttt{\textbackslash eps}), with extra wrapper packages (bm, bbm, dsfont, xspace, cancel, mathtools). The delivered numbering is the unit's own. Assets: no retained span contains a figure environment, a graphics-inclusion command, a PDF-page inclusion, a table or a TikZ picture; no image file is copied into this package, the package declares no asset dependency and no image file is redistributed with it. Licence: version arXiv:2607.20960v1 of this article is distributed under the Creative Commons Attribution 4.0 International licence, as recorded in the arXiv licence metadata for this exact version and shown on the abstract page https://arxiv.org/abs/2607.20960v1. A full rights notice is delivered at licenses/arxiv-2607.20960-CC-BY-4.0.txt. Characters: every retained excerpt is pure ASCII, so the delivered engine, XeLaTeX with its default Latin Modern text font, reports no missing character.
\begin{align}
 N_H(x)&=x\log x+(2\gamma-1)x+\Delta_H(x),
      \label{eq:def-hyperbola}\\
 N_C(x)&=\pi x+\Delta_C(x).
      \label{eq:def-circle}
\end{align}

\begin{equation}\label{eq:odd-main}
 \FF(n)=\frac{\pi}{8}n
 +\frac14\left(\Delta_C(n)-\Delta_C\left(\frac n2\right)\right)
 +O(n^\eps).
\end{equation}

\begin{equation}\label{eq:even-main}
 \FF(n)=\frac{3\log 2}{4}n
 +2\Delta_H\left(\frac n2\right)
 -5\Delta_H\left(\frac n4\right)
 +2\Delta_H\left(\frac n8\right)
 +O(n^\eps).
\end{equation}

\begin{proposition}\label{prop:converse}
Fix a real number $\theta$ with $0\leq\theta<1$.
Suppose that, for every $\eps>0$,
\[
 \FF(n)-\frac{\pi}{8}n=O(n^{\theta+\eps})
 \qquad (n\text{ odd}).
\]
Then $\Delta_C(x)=O(x^{\theta+\eps})$ for every $\eps>0$.
Suppose instead that, for every $\eps>0$,
\[
 \FF(n)-\frac{3\log 2}{4}n=O(n^{\theta+\eps})
 \qquad (n\text{ even}).
\]
Then $\Delta_H(x)=O(x^{\theta+\eps})$ for every $\eps>0$.
\end{proposition}

\begin{proof}
Assume first the bound along odd indices. Formula \eqref{eq:odd-main}, applied
with $n=2m+1$, gives
\[
 \Delta_C(2m+1)-\Delta_C\left(m+\frac12\right)
 =O(m^{\theta+\eps}).
\]
Since a sum of two integer squares is an integer,
$N_C(m+1/2)=N_C(m)$. Hence
$\Delta_C(m+1/2)=\Delta_C(m)-\pi/2$. Jacobi's identity and the
divisor bound also give
\[
 \Delta_C(j)-\Delta_C(j-1)=r_2(j)-\pi=O(j^\eps).
\]
These estimates yield, for every integer $j\geq 2$,
\[
 \Delta_C(j)=\Delta_C\left(\left\lfloor\frac j2\right\rfloor\right)
 +O(j^{\theta+\eps}).
\]
Iteration gives $\Delta_C(j)=O(j^{\theta+\eps})$. The same estimate holds
for real $x$, since $N_C(x)=N_C(\lfloor x\rfloor)$.

Assume next the bound along even indices. Formula \eqref{eq:even-main}, with
$n=8m$, gives
\[
 2\Delta_H(4m)-5\Delta_H(2m)+2\Delta_H(m)
 =O(m^{\theta+\eps}).
\]
For every positive integer $m$, set
\[
 B(m)=\Delta_H(2m)-2\Delta_H(m).
\]
In terms of $B$, the preceding bound is
\[
 2B(2m)-B(m)=O(m^{\theta+\eps}).
\]
The divisor bound and \eqref{eq:def-hyperbola} imply
\[
 \Delta_H(m+1)-\Delta_H(m)=O(m^\eps),
\]
so that $B(m+1)-B(m)=O(m^\eps)$. The two bounds for $B$ combine to give
\[
 B(j)=\frac12B\left(\left\lfloor\frac j2\right\rfloor\right)
 +O(j^{\theta+\eps})
\]
for every integer $j\geq 2$. Binary iteration shows that
\begin{equation}\label{eq:B-bound}
 B(j)=O(j^{\theta+\eps}).
\end{equation}

Fix $\eps>0$, and choose $\eta$ with
$0<\eta<\eps$ and $\theta+\eta<1$. For every positive integer $J$, the
definition of $B$ gives
\[
 \frac{\Delta_H(m)}m
 =\frac{\Delta_H(2^Jm)}{2^Jm}
 -\sum_{j=0}^{J-1}\frac{B(2^jm)}{2^{j+1}m}.
\]
The first term tends to zero as $J$ tends to infinity, since the known bounds
for the divisor error term imply $\Delta_H(x)=o(x)$. Estimate
\eqref{eq:B-bound}, used with the exponent $\theta+\eta$, makes the series
convergent and gives
\[
 \Delta_H(m)=O(m^{\theta+\eta})=O(m^{\theta+\eps}).
\]
This proves the required bound at integer arguments. Finally,
$N_H(x)=N_H(\lfloor x\rfloor)$, and the remaining main term varies by
$O(\log x)$ on an interval of length $1$. This extends the estimate from
integers to real $x$.
\end{proof}

\end{document}
